\documentclass[10pt,a4paper]{amsart}
\usepackage[margin=19mm]{geometry}
\usepackage{amsmath,amssymb,mathtools}
\usepackage[scr=rsfs,cal=euler]{mathalfa}
\usepackage{xfrac}
\usepackage[unicode,hidelinks,bookmarks=false]{hyperref}
\newcommand{\ie}{i.e.\ }
\newcommand{\cf}{cf.\ }
\newcommand{\resp}{respectively\ }
\newcommand{\Subsection}{\subsection}
\providecommand{\nobreakdash}{\nobreak\hbox{-}}
\setlength{\emergencystretch}{2em}
\multlinegap0pt
\usepackage{bm}
\usepackage{bbm}
\usepackage{dsfont}
\usepackage{xspace}
\usepackage{cancel}
\usepackage{mathtools}
\usepackage{amsthm}
\makeatletter
\@ifundefined{corollary}{\newtheorem{corollary}{Corollary}}{}
\makeatother
\title[Effect of sampling on the descriptive distributions of corre]{Effect of sampling on the descriptive distributions of correlated random walks}
\author{Joseph D. Bailey, Jessica Claridge}
\date{Source: arXiv:2607.03186v1 (CC BY 4.0)}
\begin{document}
\maketitle

\noindent\emph{Source and licence.} Effect of sampling on the descriptive distributions of correlated random walks. Version arXiv:2607.03186v1 (CC BY 4.0), distributed under the Creative Commons Attribution 4.0 International licence, as recorded in the arXiv licence metadata for this exact version and shown on the abstract page https://arxiv.org/abs/2607.03186v1. Full rights notice: licenses/arxiv-2607.03186-CC-BY-4.0.txt.\par\smallskip

\noindent\textit{Editorial scope.} Counted unit: the Corollary whose source label is \texttt{None} in the article (source0.tex lines 216--220, source label \texttt{None}), delivered together with the article's own complete original proof of it (source0.tex lines 221--244, one \texttt{proof} environment of 24 lines). No mathematics is written, repaired or re-derived: statement and proof are the article's own text line for line. Required context retained uncounted: none. Every definition, display and label of those lines is retained unchanged. Omitted from this package and not redistributed: 1--215, 245--941. Nothing inside a retained excerpt is omitted or altered: every retained body is delivered exactly as the article's lines are written, so no omission receipt of any kind accompanies this package. Citations: the retained excerpts carry no citation command at all, so no bibliography entry of the article is transcribed and no reference is redistributed. Reference adaptation: none was needed and none was made; every reference inside the delivered unit resolves inside it, so no unresolved reference or citation is produced. Printed slips kept literal: no printed word, symbol, hypothesis or number of the counted statement or of its proof was altered. Layout and class: the article's own class is \texttt{article} with packages including geometry, authblk, fontenc, inputenc, babel, graphicx, float, bbm, amsmath, amssymb, amsfonts, amsthm, mathrsfs, appendix; the wrapper is the lane's pinned \texttt{amsart} editorial wrapper at 10pt on A4, which restates the theorem-like environments the retained text needs and adds the article's own macro definitions that the retained text needs (none), with extra wrapper packages (bm, bbm, dsfont, xspace, cancel, mathtools). The delivered numbering is the unit's own. Assets: no retained span contains a figure environment, a graphics-inclusion command, a PDF-page inclusion, a table or a TikZ picture; no image file is copied into this package, the package declares no asset dependency and no image file is redistributed with it. Licence: version arXiv:2607.03186v1 of this article is distributed under the Creative Commons Attribution 4.0 International licence, as recorded in the arXiv licence metadata for this exact version and shown on the abstract page https://arxiv.org/abs/2607.03186v1. A full rights notice is delivered at licenses/arxiv-2607.03186-CC-BY-4.0.txt. Characters: every retained excerpt is pure ASCII, so the delivered engine, XeLaTeX with its default Latin Modern text font, reports no missing character.
\begin{corollary}
    For a correlated random walk with constant step-length distribution, that is, for all $\ell$ we have $\Lambda(\ell)=\ell^*\in \mathbb{R}$, and turning angles, $\theta_t$, drawn from a circular distribution $f_\circ$ centred at 0, the distribution of step-lengths of the random walk formed by sub-sampling at rate $r=2$ is given by \begin{equation}
        \Lambda^{\{2\}}(\ell^{\{2\}})=\frac{2}{\sqrt{(\ell^*)^2-\left(\frac{\ell^{\{2\}}}{2}\right)^2}}f_\circ\left(2\arccos{\left(\frac{\ell^{\{2\}}}{2\ell^*}\right)}\right)
    \end{equation} which has range $[0,2\ell^*]$.  
\end{corollary}

\begin{proof}
Consider the consecutive points of a CRW $X_{t-1},X_t,X_{t+1}$.  We have $|X_{t-1}X_t|=|X_tX_{t+1}|=\ell^*$ by the properties of $\Lambda$, we are interested in describing the distribution of $\ell_t^{\{2\}}=|X_{t-1}X_{t+1}|$.  Using the cosine rule, The length of the vector between points $X_{t-1}$ and $X_{t+1}$ for all $t$ is simply
   \begin{align}\nonumber
     \lvert X_{t-1}X_{t+1}\rvert^2&=\lvert X_{t-1}X_t\rvert^2+\lvert X_{t}X_{t+1}\rvert^2-2\lvert X_{t-1}X_t\rvert\lvert X_{t}X_{t+1}\rvert\cos{\left(\pi-\theta_{X_t}\right)}\\\nonumber
       \left(\ell_t^{\{2\}}\right)^2&=2\ell^{*2}+2\ell^{*2}\cos{\left(\theta_{X_t}\right)}\\
       %\nonumber &=2\ell^2(1+\cos{\left(\theta_{X_t}\right)})\\
       \Rightarrow \ell_t^{\{2\}}&=\sqrt{2\ell^{*2}(1+\cos{\left(\theta_{X_t}\right)})}=2\ell^*\cos{\frac{\theta_{X_t}}{2}}
       \label{eqn: 12}
   \end{align}
This shows that the distribution of $\ell^{\{2\}}$ is given in terms of the distribution of $\theta$, with $\ell_t^{\{2\}}$ taking values in $[0,2\ell^*]$ as $\theta_{X_t}\in(-\pi,\pi]$.  Now, as all $\theta_{X_t}$ are i.i.d, to find the distribution, $\Lambda^{\{2\}}(\ell^{\{2\}})$, we have
   \begin{align}\nonumber
       F_{\Lambda^{\{2\}}}(\ell^{\{2\}})&=P(L\leq \ell^{\{2\}})\\\nonumber
       &=P\left(2\ell^*\cos{\frac{\theta}{2}}\leq \ell^{\{2\}}\right)\\\nonumber
       &=2P\left(\theta\leq 2\arccos{\left(\frac{\ell^{\{2\}})}{2\ell^*}\right)}\right)\\
       &=2F_\theta\left(2\arccos{\left(\frac{\ell^{\{2\}})}{2\ell^*}\right)}\right)
   \end{align}
   where $F_{\Lambda^{\{2\}}}, F_\theta$ are the CDFs of $\Lambda^{\{2\}}$ and $f_\circ$ respectively.  The coefficient of 2 appears due to $\theta$ having domain $[-\pi,\pi]$, arccos having the restricted range of $[0,\pi]$ and cosine being an even function.  Hence, the distribution $\Lambda^{\{2\}}$ is given by
   \begin{align}\nonumber
       \Lambda^{\{2\}}(\ell^{\{2\}})=\frac{dF_{\Lambda^{\{2\}}}(\ell^{\{2\}})}{d\ell^{\{2\}}}&=2\frac{dF_\theta}{d\theta}\Bigg\lvert\frac{d\theta}{d\ell^{\{2\}}}\Bigg\rvert\\\nonumber
       &=2\frac{dF_\theta}{d\theta}\Bigg\lvert\frac{d}{d\ell^{\{2\}}}\left(2\arccos{\left(\frac{\ell^{\{2\}}}{2\ell^*}\right)}\right)\Bigg\rvert\\
       &=\frac{2}{\sqrt{\ell^{*2}-\left(\frac{\ell^{\{2\}}}{2}\right)^2}}f_\circ\left(2\arccos{\left(\frac{\ell^{\{2\}}}{2\ell^*}\right)}\right),
   \end{align}
Noting that here we have $\frac{dF_\theta}{d\theta}=f_\circ(\theta)$ and $\theta=2\arccos{\left(\frac{\ell^{\{2\}})}{2\ell^*}\right)}$.
\end{proof}

\end{document}
